% Workbench module: exact S^8 constraints for the higher-rung search.
% Date: 2026-09-03.  No standalone preamble is included.
% All labels in this module use the prefix hrs8:.

\section{Exact constraints on an \texorpdfstring{$S^8$}{S8} higher-rung analogue}
\label{hrs8:sec:constraints}

The phrase ``an $S^8$ analogue'' can denote three mathematically different
proposals, which must not be conflated.  A literal smooth complex-fourfold
analogue is impossible because a homotopy $8$-sphere is not almost complex.
A merely real map to $S^2$ with generic fibre $T^6$ is not excluded if it has
genuine singular fibres, but its local fillings must satisfy the exact Euler,
fundamental-group, Wang, and Mayer--Vietoris tests below.  Finally, explicit
real stratified quotient models with boundary or total space $S^8$ do exist;
they do not thereby acquire a complex-fourfold structure.

There is also a different, well-typed alternative: retain the four-real-
dimensional torus fibres of the reconstructed $S^6$ candidate and place
$SU(3)$ bundle data on its regular locus.  Subsection
\ref{hrs8:subsec:su3-regular-locus} computes the exact extension obstruction.
No gauge-theoretic or physical conclusion is asserted without a specified
bundle, connection, trace convention, boundary condition, and analytic
operator statement.

\subsection{The literal complex-fourfold proposal}
\label{hrs8:subsec:no-almost-complex}

\subsubsection{Characteristic-class obstruction}
\label{hrs8:subsubsec:characteristic-class}

Let $M$ be a smooth manifold homotopy equivalent to $S^8$, and suppose that
$J\colon TM\to TM$ satisfies $J^2=-1$.  Then
\[
 E=(TM,J)
\]
is a complex vector bundle of rank four whose underlying oriented real bundle
is $TM$.  Since
\[
 H^2(M;\mathbb Z)=H^4(M;\mathbb Z)=H^6(M;\mathbb Z)=0,
\]
one has
\[
 c_1(E)=c_2(E)=c_3(E)=0.
\]
Orient $M$ by $J$, and choose $u\in H^8(M;\mathbb Z)$ with
$\langle u,[M]\rangle=1$.  The top-Chern--Euler identity and the
Poincar\'e--Hopf theorem give
\begin{equation}
 c_4(E)=e(TM)=\chi(M)u=2u.
 \label{hrs8:eq:top-chern-euler}
\end{equation}
For a complex bundle, the second Pontryagin class of the underlying real
bundle is
\[
 p_2(E_{\mathbb R})
 =c_2(E)^2-2c_1(E)c_3(E)+2c_4(E).
\]
Consequently,
\begin{equation}
 p_2(TM)=4u.
 \label{hrs8:eq:p2-complex}
\end{equation}
On the other hand, $H^4(M;\mathbb Z)=0$ gives $p_1(TM)=0$, and the middle
cohomology vanishes, so $\operatorname{sign}(M)=0$.  Hirzebruch's
eight-dimensional signature formula is
\[
 \operatorname{sign}(M)
 =\left\langle\frac{7p_2(TM)-p_1(TM)^2}{45},[M]\right\rangle.
\]
It forces $\langle p_2(TM),[M]\rangle=0$, contradicting
\eqref{hrs8:eq:p2-complex}.  For the standard embedded sphere, the signature
step may instead be replaced by the explicit stable trivialisation
\[
 TS^8\oplus\varepsilon^1_{\mathbb R}\cong\varepsilon^9_{\mathbb R},
\]
which gives $p_2(TS^8)=0$ directly.  Hence
\begin{equation}
 \boxed{S^8\text{ admits no almost-complex structure}.}
 \label{hrs8:eq:no-almost-complex}
\end{equation}
Integrability is stronger than being almost complex, so a complex-fourfold
structure is excluded a fortiori.

\subsubsection{Independent arithmetic checks}
\label{hrs8:subsubsec:arithmetic-checks}

Borel and Serre prove that $S^{2n}$ admits no almost-complex structure for
$n\geq4$, including arbitrary differentiable structures on the sphere.  In
the $S^8$ specialization, their mod-$3$ reduced-power identity is
\[
 \mathcal P^1(c_2)
 =c_1^2c_2-2c_2^2-c_1c_3+4c_4.
\]
All lower Chern classes vanish on $S^8$, so it would follow that
$c_4=0\pmod 3$, whereas \eqref{hrs8:eq:top-chern-euler} gives
$c_4=\pm2u\neq0\pmod3$.

Bott integrality gives a second independent check.  For a complex bundle
over $S^{2n}$, the top Chern number is divisible by $(n-1)!$.  At $n=4$,
\[
 6\mid\langle c_4(E),[S^8]\rangle,
\]
which is incompatible with the Euler number $\pm2$.

\subsubsection{Exact scope of the obstruction}
\label{hrs8:subsubsec:scope}

The obstruction rules out a smooth complex fourfold diffeomorphic to $S^8$;
a proper holomorphic family of complex three-tori over $\mathbb P^1$ whose
smooth total space is diffeomorphic to $S^8$; and an almost-complex structure
on all of $TS^8$, whether or not integrable.  It also rules out a symplectic
or K\"ahler structure: if $\omega$ were symplectic, then
$H^2(S^8;\mathbb R)=0$ would give $\omega=d\alpha$, while
\[
 \int_{S^8}\omega^4
 =\int_{S^8}d(\alpha\wedge\omega^3)=0,
\]
contradicting symplectic volume.

It does not rule out a complex structure on an open subset, a stable complex
structure, a singular complex analytic space whose tangent sheaf is not a
vector bundle, or a real stratified quotient whose underlying boundary is
$S^8$.  The standard sphere is stably complex because
\[
 TS^8\oplus\varepsilon^2_{\mathbb R}\cong\varepsilon^{10}_{\mathbb R},
\]
but this is not a complex structure on the actual rank-eight tangent bundle.
Moreover, the usual link of an isolated complex singularity of complex
dimension $d$ has real dimension $2d-1$.  Thus an $S^8$ obtained as a quotient
boundary is not, merely for that reason, an ordinary complex-analytic link.

\subsection{Requirements for a real Fabel-style \texorpdfstring{$S^8$}{S8} fibration}
\label{hrs8:subsec:real-fibration}

The reconstructed $S^6$ candidate has a map
\[
 f\colon X^6\longrightarrow B\cong\mathbb P^1
\]
whose regular fibre is a complex two-torus, hence a real $T^4$, and whose
three special fibres comprise one non-normal toric degeneration and two
multiple bielliptic fibres.  The direct dimension analogue on a real
eight-manifold would be a map
\[
 g\colon M^8\longrightarrow S^2
\]
with regular fibre $T^6$.  If $M$ is required to be a smooth complex
fourfold and the regular fibres are complex three-tori, the obstruction in
Subsection \ref{hrs8:subsec:no-almost-complex} ends the proposal.  If $g$ is
only a real smooth or stratified map, the proposal remains topologically
meaningful.

The adjective ``singular'' is essential.  There is no fibre bundle
\[
 T^6\longrightarrow S^8\longrightarrow S^2.
\]
Indeed, its long exact homotopy sequence would contain
\[
 \pi_2(S^8)=0\longrightarrow\pi_2(S^2)=\mathbb Z
 \xrightarrow{\partial}\pi_1(T^6)=\mathbb Z^6
 \longrightarrow\pi_1(S^8)=0.
\]
Exactness would require $\partial\colon\mathbb Z\to\mathbb Z^6$ to be
surjective, which is impossible.  The Euler identity for a fibre bundle gives
the same conclusion:
\[
 \chi(S^8)=2\neq\chi(S^2)\chi(T^6)=2\cdot0=0.
\]

\subsection{Euler constraint for a finite multipole model}
\label{hrs8:subsec:euler}

Let $g\colon M^8\to S^2$ be a proper triangulable map with finite
critical-value set $\Sigma=\{p_1,\ldots,p_s\}$, and suppose that every fibre
over $S^2\setminus\Sigma$ is $T^6$.  Euler integration for a constructible
map gives
\begin{align}
 \chi(M)
 &=\chi(S^2\setminus\Sigma)\chi(T^6)
   +\sum_{i=1}^s\chi\bigl(g^{-1}(p_i)\bigr)\notag\\
 &=\sum_{i=1}^s\chi(F_i).
 \label{hrs8:eq:euler-integration}
\end{align}
Equivalently, one may cut out small discs, use multiplicativity on the regular
bundle, and note that each boundary mapping torus has Euler characteristic
zero.  For $M\cong S^8$ one therefore needs
\begin{equation}
 \boxed{\sum_{i=1}^s\chi(F_i)=2.}
 \label{hrs8:eq:euler-fibres}
\end{equation}
Multiple fibres that are finite free quotients of tori have Euler
characteristic zero and cannot account for $S^8$.  At least one special fibre
must contain a genuine collapse or fixed stratum with nonzero Euler
characteristic.  The numerical pattern consisting of two Euler-zero multiple
fibres and one fibre of Euler characteristic two is compatible with
\eqref{hrs8:eq:euler-fibres}, but that numerical compatibility is not a
construction.

\subsection{Fundamental group and monodromy}
\label{hrs8:subsec:pi1}

Put $B^\circ=S^2\setminus\Sigma$, $E=g^{-1}(B^\circ)$, and
$\Lambda=H_1(T^6;\mathbb Z)\cong\mathbb Z^6$.  The regular bundle has
monodromy
\[
 \rho\colon\pi_1(B^\circ)\longrightarrow\operatorname{GL}(\Lambda).
\]
If the fibres and total space are oriented compatibly, its image lies in
$\operatorname{SL}(6,\mathbb Z)$.  For oriented peripheral loops
$\gamma_1,\ldots,\gamma_s$ satisfying $\gamma_1\cdots\gamma_s=1$, the
matrices $M_i=\rho(\gamma_i)$ satisfy
\begin{equation}
 M_1M_2\cdots M_s=I.
 \label{hrs8:eq:monodromy-product}
\end{equation}

When $s\geq1$, $B^\circ$ has $\pi_2(B^\circ)=0$.  The homotopy sequence of
the regular bundle therefore gives the exact extension
\begin{equation}
 1\longrightarrow\Lambda\longrightarrow\pi_1(E)
 \longrightarrow\pi_1(B^\circ)\longrightarrow1.
 \label{hrs8:eq:regular-pi1-extension}
\end{equation}
After choosing a section and a compatible basepoint system,
\eqref{hrs8:eq:regular-pi1-extension} becomes an affine semidirect
presentation; without those choices it remains the displayed exact
extension.

Let $N_i=g^{-1}(D_i)$ be the filling over a small closed disc $D_i$ about
$p_i$, and let $A_i=E\cap N_i$, a $T^6$ mapping torus.  Van Kampen gives
\begin{equation}
 \pi_1(M)\cong
 \pi_1(E)\Big/
 \Big\langle\!\Big\langle
  \ker\bigl(\pi_1(A_i)\to\pi_1(N_i)\bigr),\ 1\leq i\leq s
 \Big\rangle\!\Big\rangle.
 \label{hrs8:eq:van-kampen-fillings}
\end{equation}
Thus an $S^8$ construction must derive the actual local inclusion maps and
prove that their kernels normally generate all of
\eqref{hrs8:eq:regular-pi1-extension}.  Monodromy matrices or Smith forms
alone do not prove that \eqref{hrs8:eq:van-kampen-fillings} is trivial.

\subsection{Integral homology: the exact gluing tests}
\label{hrs8:subsec:homology}

Replace $B^\circ$ by a compact sphere with $s$ open discs removed and replace
$E,N_i$ by compact collar-compatible pieces.  Define the actual
Mayer--Vietoris maps
\begin{equation}
 \alpha_q=
 \bigoplus_{i=1}^s\bigl((j_{E,i})_*,-(j_{N,i})_*\bigr)\colon
 \bigoplus_iH_q(A_i;\mathbb Z)\longrightarrow
 H_q(E;\mathbb Z)\oplus\bigoplus_iH_q(N_i;\mathbb Z).
 \label{hrs8:eq:mv-map}
\end{equation}
Because the pieces have nonempty boundary, their top homology vanishes.  The
Mayer--Vietoris sequence and
$H_q(S^8;\mathbb Z)=0$ for $1\leq q\leq7$ imply the following necessary and
sufficient conditions:
\begin{equation}
\begin{array}{c|l}
q&\text{condition on }\alpha_q\\ \hline
0&\alpha_0\text{ is injective and }\operatorname{coker}\alpha_0\cong\mathbb Z,\\
1\leq q\leq6&\alpha_q\text{ is an isomorphism},\\
7&\alpha_7\text{ is surjective and }\ker\alpha_7\cong\mathbb Z.
\end{array}
\label{hrs8:eq:mv-tests}
\end{equation}
The generator of $\ker\alpha_7$ is the boundary of the fundamental class of
the closed eight-manifold obtained by gluing the pieces.  Equations
\eqref{hrs8:eq:mv-map}--\eqref{hrs8:eq:mv-tests}, using the actual inclusion
maps, are the integral replication test for a proposed $S^8$ multipole
construction.

Each boundary group is constrained by monodromy.  If $A_i$ is the mapping
torus of $\phi_i\colon T^6\to T^6$ and
$M_i=(\phi_i)_*|_{H_1(T^6;\mathbb Z)}$, the Wang sequence gives
\begin{equation}
 0\longrightarrow
 \operatorname{coker}\bigl(I-\textstyle\bigwedge^qM_i\bigr)
 \longrightarrow H_q(A_i;\mathbb Z)
 \longrightarrow
 \ker\bigl(I-\textstyle\bigwedge^{q-1}M_i\bigr)
 \longrightarrow0.
 \label{hrs8:eq:wang}
\end{equation}
The right group is free, so \eqref{hrs8:eq:wang} splits as an abelian group,
although not canonically.  Since the regular base retracts to a wedge of
$s-1$ circles, the homological Serre spectral sequence gives the filtration
\begin{equation}
 0\longrightarrow
 H_0\!\left(B^\circ;\textstyle\bigwedge^q\Lambda\right)
 \xrightarrow{\ \iota_q\ } H_q(E;\mathbb Z)
 \xrightarrow{\ \pi_q\ }
 H_1\!\left(B^\circ;\textstyle\bigwedge^{q-1}\Lambda\right)
 \longrightarrow0.
 \label{hrs8:eq:serre-regular-part}
\end{equation}
To type both edge terms explicitly, put
\(V_q=\bigwedge^q\Lambda\), choose oriented free generators
\(\gamma_1,\ldots,\gamma_{s-1}\) of \(\pi_1(B^\circ)\), and write
\(M_j=\rho(\gamma_j)\).  In the Wang sign convention the cellular differential is
\begin{equation}
 d_q\colon\bigoplus_{j=1}^{s-1}V_q\longrightarrow V_q,\qquad
 (v_j)_j\longmapsto
 \sum_{j=1}^{s-1}\left(I-\textstyle\bigwedge^qM_j\right)v_j,
 \label{hrs8:eq:serre-cellular-differential}
\end{equation}
so the two homological edge groups are
\begin{equation}
 H_0(B^\circ;V_q)
 =\operatorname{coker}d_q
 =(V_q)_\rho
 =\frac{V_q}{
   \sum_{j=1}^{s-1}(I-\bigwedge^qM_j)V_q},
 \qquad
 H_1(B^\circ;V_{q-1})=\ker d_{q-1}.
 \label{hrs8:eq:serre-coinvariants-kernel}
\end{equation}
Thus the coinvariants are the Serre filtration subobject and the
base-one-cycle group is the quotient.  Since \(\ker d_{q-1}\) is a subgroup
of a free abelian group, it is free, and
\eqref{hrs8:eq:serre-regular-part} splits noncanonically as an abstract
abelian group.

Dually, set \(V_q^\vee=\bigwedge^q\Lambda^\vee\); its monodromy along
\(\gamma_j\) is \(\bigwedge^q(M_j^{-\mathsf T})\).  The cohomological
filtration is
\begin{equation}
 0\longrightarrow H^1(B^\circ;V_{q-1}^\vee)
 \longrightarrow H^q(E;\mathbb Z)
 \longrightarrow H^0(B^\circ;V_q^\vee)\longrightarrow0,
 \qquad
 \begin{aligned}
 H^0(B^\circ;V_q^\vee)
   &=(V_q^\vee)^\rho\\
   &=\bigcap_{j=1}^{s-1}
     \ker\!\left(I-\textstyle\bigwedge^q(M_j^{-\mathsf T})\right).
 \end{aligned}
 \label{hrs8:eq:serre-cohomology-invariants}
\end{equation}
Accordingly, each boundary inclusion in \(\alpha_q\) must send its Wang
coinvariant filtration term through \(\iota_q\), and its base-one-cycle
term must be tested after \(\pi_q\); neither exact sequence supplies a
canonical splitting.  The corrected sequence identifies \(H_q(E;\mathbb Z)\)
noncanonically with the direct sum of its two edge groups, so a rank or
torsion count made from that direct sum survives.  The natural edge maps do
not: every map-level Mayer--Vietoris test must use
\eqref{hrs8:eq:serre-regular-part} in the displayed order.  Equations
\eqref{hrs8:eq:mv-map}--\eqref{hrs8:eq:serre-cohomology-invariants} retain
every lattice, monodromy, kernel, coinvariant, invariant, and inclusion
that a correct proof must compute.

\subsection{The torus-action bound and a maximal real orbit model}
\label{hrs8:subsec:torus-action}

For a smooth torus action on a compact manifold,
\[
 \chi(M)=\chi(M^T).
\]
Since $\chi(S^8)=2$, every torus action on $S^8$ has a fixed point.  Equip
$S^8$ with a $T$-invariant Riemannian metric.  At a fixed point the isotropy
representation is a sum of four real two-dimensional weight planes and
possibly trivial summands.  If the global action is effective, this tangent
representation is faithful: an isometry fixing a point with identity
derivative is the identity on the connected manifold.  Four weights span a
character lattice of rank at most four.  Hence
\begin{equation}
 \boxed{\operatorname{rank}T\leq4.}
 \label{hrs8:eq:torus-rank}
\end{equation}
In particular, a generic $T^6$ fibration of $S^8$ cannot be the orbit map of
a global effective torus action.  This does not rule out a monodromic
$T^6$-bundle over the regular part of $S^2$, because such a bundle need not
supply a global $T^6$ action.

The bound is sharp.  Write
\[
 S^8=\left\{(z_1,z_2,z_3,z_4,y)\in\mathbb C^4\oplus\mathbb R:
       \sum_{i=1}^4|z_i|^2+y^2=1\right\},
\]
and let $T^4$ rotate the four complex coordinates.  The orbit map is
\[
 q(z_1,z_2,z_3,z_4,y)
 =(|z_1|,|z_2|,|z_3|,|z_4|,y),
\]
with image
\[
 Q=\left\{(r_1,r_2,r_3,r_4,y):r_i\geq0,
       \sum r_i^2+y^2=1\right\}.
\]
Two points have the same image exactly when they lie in the same orbit, so
compact--Hausdorff descent gives $S^8/T^4\cong Q$.  Away from the two poles,
write $r_i=\sqrt{1-y^2}\,u_i$.  Then
\[
 (u_1,u_2,u_3,u_4)\in
 \{u_i\geq0,\ \sum u_i^2=1\}\cong\Delta^3,
\]
and the values $y=\pm1$ collapse the two ends.  Therefore
\begin{equation}
 \boxed{S^8/T^4\cong\operatorname{Susp}(\Delta^3)\cong B^4.}
 \label{hrs8:eq:t4-quotient}
\end{equation}
The locus on which all $r_i>0$ has free $T^4$ orbit.  Exactly $j$ zero
complex coordinates give isotropy $T^j$, and the two poles are fixed by the
whole torus.  This is a rigorous two-pole singular torus compactification of
$S^8$, but it has four-dimensional quotient and generic $T^4$ orbit.  It is
not a three-special-fibre map to $S^2$.

\subsection{A three-column Gram quotient whose quotient boundary is \texorpdfstring{$S^8$}{S8}}
\label{hrs8:subsec:gram}

Let
\[
 V=\operatorname{Mat}_{3\times3}(\mathbb C),
 \qquad \dim_{\mathbb R}V=18,
\]
and let $U(3)$ act on the left.  Define
\begin{equation}
 \Gamma\colon V\longrightarrow\operatorname{Herm}_3(\mathbb C)_{\geq0},
 \qquad \Gamma(A)=A^*A.
 \label{hrs8:eq:gram-map}
\end{equation}
The map is surjective because every $Q\geq0$ is the Gram matrix of
$Q^{1/2}$.  It separates $U(3)$-orbits.  Indeed, if $A^*A=B^*B$, the rule
$Ax\mapsto Bx$ is a well-defined isometry from $\operatorname{im}A$ to
$\operatorname{im}B$: equality of the Gram matrices implies both equality of
kernels and preservation of the Hermitian inner product.  Extend this
isometry to a unitary $U$ of $\mathbb C^3$; then $B=UA$.  Hence
\begin{equation}
 V/U(3)\cong\operatorname{Herm}_3(\mathbb C)_{\geq0}.
 \label{hrs8:eq:gram-orbit-space}
\end{equation}

For the Frobenius unit ball
$\mathbb B(V)=\{A:\operatorname{tr}(A^*A)\leq1\}$,
\eqref{hrs8:eq:gram-map} gives
\begin{equation}
 \mathbb B(V)/U(3)
 \cong K:=\{Q\in\operatorname{Herm}_3(\mathbb C):Q\geq0,
             \operatorname{tr}Q\leq1\}.
 \label{hrs8:eq:gram-ball-quotient}
\end{equation}
The real vector space of Hermitian $3\times3$ matrices has dimension nine.
The set $K$ is compact, convex, and has nonempty interior
$\{Q>0,\operatorname{tr}Q<1\}$.  Radial projection from an interior point
gives a homeomorphism of pairs $(K,\partial K)\cong(B^9,S^8)$, and therefore
\begin{equation}
 \boxed{\partial\bigl(\mathbb B^{18}/U(3)\bigr)\cong S^8.}
 \label{hrs8:eq:gram-boundary}
\end{equation}

Two qualifications are essential.  First, the quotient of the source sphere
is not the boundary in \eqref{hrs8:eq:gram-boundary}:
\[
 S^{17}/U(3)\cong\{Q\geq0:\operatorname{tr}Q=1\}.
\]
The space on the right is an eight-dimensional compact convex body and hence
is homeomorphic to $B^8$.  The rest of $\partial K$ consists of
rank-deficient Gram matrices with $\operatorname{tr}Q<1$.  Thus taking the
boundary and taking the orbit quotient do not commute in this construction.

Second, strict $SU(3)$ is different.  For invertible matrices, $A^*A$ does
not record the phase of $\det A$, while left multiplication by $SU(3)$
preserves $\det A$.  The strict quotient is described by
\begin{equation}
 (Q,d)=(A^*A,\det A),
 \qquad Q\geq0,
 \qquad |d|^2=\det Q.
 \label{hrs8:eq:strict-su3-quotient}
\end{equation}
It has generic real dimension ten, not nine.  The construction in
\eqref{hrs8:eq:gram-boundary} therefore uses
\[
 U(3)\cong(SU(3)\times U(1))/\mu_3,
\]
not colour $SU(3)$ alone.  It is a precise $S^8$ real quotient boundary, but
no morphism has been produced from its Gram or rank strata to the period
matrix, torus bundle, logarithmic twists, or three special fibres of the
Fabel--Alp\"oge construction.

\subsection{Retaining the \texorpdfstring{$T^4/\mathbb R^4$}{T4/R4} fibre geometry and placing \texorpdfstring{$SU(3)$}{SU(3)} on the regular locus}
\label{hrs8:subsec:su3-regular-locus}

This subsection classifies the required topological bundle data and its
obstruction to crossing the singular fibres.  It makes no physical
identification.

\subsubsection{The regular locus and its fourth cohomology}
\label{hrs8:subsubsec:regular-cohomology}

Let $X\cong_{\mathrm{diff}}S^6$ be the reconstructed candidate and put
\[
 D=W\sqcup S_1\sqcup S_2,
 \qquad U=X\setminus D=f^{-1}(B^\circ),
 \qquad B^\circ=\mathbb P^1\setminus\{p_0,p_1,p_2\}.
\]
Here $W$ is the non-normal central fibre and $S_1,S_2$ are the underlying
smooth bielliptic surfaces of the two multiple fibres.  The reconstructed
integral calculations give
\[
 H_1(W;\mathbb Z)\cong\mathbb Z^2,
 \qquad H_1(S_1;\mathbb Z)\cong H_1(S_2;\mathbb Z)\cong\mathbb Z^2.
\]
The compact analytic set $D$ is triangulable, so Alexander duality in
$X\cong S^6$ applies:
\[
 \widetilde H^i(U;\mathbb Z)
 \cong\widetilde H_{5-i}(D;\mathbb Z).
\]
At $i=4$, it gives
\begin{equation}
 \boxed{H^4(U;\mathbb Z)\cong H_1(D;\mathbb Z)\cong\mathbb Z^6.}
 \label{hrs8:eq:h4-regular-locus}
\end{equation}

There is an independent fibrewise description.  The pair of pants
$B^\circ$ retracts to a wedge of two circles, and the $T^4$ bundle retracts
to its restriction over that graph.  Hence $U$ has the homotopy type of a CW
complex of dimension at most five.  If $\mathcal V^q=R^qf_*\mathbb Z$ is the
monodromy local system, the Serre spectral sequence has only columns zero and
one and yields
\begin{equation}
 0\longrightarrow H^1(B^\circ;\mathcal V^3)
 \longrightarrow H^4(U;\mathbb Z)
 \xrightarrow{\operatorname{res}_F}
 H^0(B^\circ;\mathcal V^4)
 \longrightarrow0.
 \label{hrs8:eq:h4-serre}
\end{equation}
Holomorphic fibre monodromy preserves the real orientation, so
$H^0(B^\circ;\mathcal V^4)\cong\mathbb Z$.  Thus one coordinate of
\eqref{hrs8:eq:h4-regular-locus} is the constant evaluation on a smooth
$T^4$ fibre; the remaining rank-five kernel in \eqref{hrs8:eq:h4-serre}
consists of mixed base--fibre classes.

\subsubsection{Exact classification of \texorpdfstring{$SU(3)$}{SU(3)} bundles on \texorpdfstring{$U$}{U}}
\label{hrs8:subsubsec:su3-classification}

Let $P\to U$ be a principal $SU(3)$ bundle and
$E_P=P\times_{SU(3)}\mathbb C^3$ its determinant-trivial rank-three complex
bundle.  Through degree five,
\[
 \pi_i(BSU(3))=0\quad(i=1,2,3,5),
 \qquad \pi_4(BSU(3))\cong\mathbb Z,
\]
with universal generator $c_2$.  Because $U$ has CW dimension at most five,
obstruction theory gives a bijection
\begin{equation}
 \boxed{
 \{\text{principal }SU(3)\text{ bundles on }U\}/\!\cong
 \ \xrightarrow{\ c_2\ }\ H^4(U;\mathbb Z)\cong\mathbb Z^6.}
 \label{hrs8:eq:su3-classification}
\end{equation}
The fibre integer
\begin{equation}
 k_F=\left\langle c_2(E_P)|_{F_b},[F_b]\right\rangle,
 \qquad b\in B^\circ,
 \label{hrs8:eq:fibre-instanton-number}
\end{equation}
is independent of $b$.  A nonzero $k_F$ is only one possible obstruction: a
class with $k_F=0$ may still have a nonzero mixed component in the kernel of
\eqref{hrs8:eq:h4-serre}.

\subsubsection{Extension across the three special fibres}
\label{hrs8:subsubsec:extension}

The cohomology sequence of the pair $(X,U)$ contains
\[
 H^4(X;\mathbb Z)\longrightarrow H^4(U;\mathbb Z)
 \xrightarrow{\delta}H^5(X,U;\mathbb Z)
 \longrightarrow H^5(X;\mathbb Z).
\]
Both outer groups vanish because $X\cong S^6$.  Therefore
\begin{equation}
 \boxed{\delta\colon H^4(U;\mathbb Z)
 \xrightarrow{\ \cong\ }H^5(X,U;\mathbb Z).}
 \label{hrs8:eq:relative-delta}
\end{equation}
Every $SU(3)$ bundle on all of $X$ has $c_2=0$, so its restriction to $U$ is
trivial by \eqref{hrs8:eq:su3-classification}.  Conversely, the trivial
bundle on $U$ extends as the trivial bundle on $X$.  Thus, for unframed
topological bundles,
\begin{equation}
 \boxed{P\to U\text{ extends across }D
 \iff c_2(P)=0
 \iff\delta c_2(P)=0.}
 \label{hrs8:eq:extension-criterion}
\end{equation}

Under Alexander duality, \eqref{hrs8:eq:relative-delta} identifies the
obstruction with a six-integer class supported on the one-cycles of
$W\sqcup S_1\sqcup S_2$.  Because the three components are disjoint and
$H_1(D)$ is their direct sum, nonzero components do not vanish merely by
declaring that charges at distinct special fibres cancel.  They cancel as an
extension obstruction only when the complete class in $\mathbb Z^6$ is zero.

This is a statement about the bundle.  Even when $c_2=0$, a specified
connection on the trivial bundle may fail to extend smoothly because of its
holonomy or growth near $D$; this is an analytic boundary-value problem not
decided by \eqref{hrs8:eq:extension-criterion}.

\subsubsection{The distinct six-dimensional clutching charge}
\label{hrs8:subsubsec:c3-clutching}

Principal $SU(3)$ bundles on the complete sphere are classified by
\[
 [S^6,BSU(3)]\cong\pi_5(SU(3))\cong\mathbb Z.
\]
If $k$ is the clutching integer, Bott integrality fixes the top Chern class as
\begin{equation}
 c_3(E_P)=\pm2k\,u\in H^6(S^6;\mathbb Z).
 \label{hrs8:eq:c3-clutching}
\end{equation}
Every such global bundle restricts to the trivial bundle on $U$, because its
restricted $c_2$ vanishes and \eqref{hrs8:eq:su3-classification} is complete.
Thus a global six-dimensional $SU(3)$ clutching charge can be concentrated in
the gluing across the deleted fibres, but it is a
$c_3/\pi_5(SU(3))$ class, not a four-dimensional $c_2$ instanton class on the
regular fibres.

Alexander duality also gives
$H^5(U;\mathbb Z)\cong\widetilde H_0(D;\mathbb Z)\cong\mathbb Z^2$.
The pair sequence therefore yields
\[
 0\longrightarrow\mathbb Z^2
 \longrightarrow H^6(X,U;\mathbb Z)
 \longrightarrow\mathbb Z
 \longrightarrow0,
\]
and hence
\begin{equation}
 H^6(X,U;\mathbb Z)\cong\mathbb Z^3.
 \label{hrs8:eq:h6-relative}
\end{equation}
After choosing the complex orientations of the three fibre components, the
three relative top classes map to the global fundamental class by the
augmentation
\begin{equation}
 (n_0,n_1,n_2)\longmapsto n_0+n_1+n_2.
 \label{hrs8:eq:augmentation}
\end{equation}
The rank-two kernel records redistribution among the three supports with the
same global total.  To promote
\eqref{hrs8:eq:h6-relative}--\eqref{hrs8:eq:augmentation} from exact
cohomological data to a relative $SU(3)$ bundle theorem with prescribed framings,
one must give the three clutching maps and compute their homotopy classes;
arbitrary relative cohomology integers are not silently declared realised.

\subsubsection{Data required before a gauge-theory conclusion}
\label{hrs8:subsubsec:gauge-data}

A mathematical colour or gauge proposal on $U$ must specify:
\begin{enumerate}
 \item a principal $SU(3)$ bundle $P\to U$, equivalently its complete class
       $c_2(P)\in\mathbb Z^6$;
 \item a connection $A$ and a fixed trace convention for its curvature
       Chern--Weil form;
 \item tubular boundary conditions or framings at each of $W,S_1,S_2$;
 \item whether the invariant under study is the fibre number
       \eqref{hrs8:eq:fibre-instanton-number}, the extension defect
       \eqref{hrs8:eq:relative-delta}, or the six-dimensional clutching
       number \eqref{hrs8:eq:c3-clutching}; and
 \item an explicit analytic statement about a Yang--Mills operator, action,
       or spectrum.
\end{enumerate}
The topology alone proves no mass gap, confinement statement, particle
identification, or Yang--Mills dynamics.  It supplies a carrier and an exact
obstruction ledger on which such a calculation can be posed.

\subsection{Consequences for the higher-rung search}
\label{hrs8:subsec:consequences}

Replacing $S^6$ by $S^8$ does not permit a literal smooth complex-fourfold
version of the Fabel--Alp\"oge construction.  A real $T^6$-over-$S^2$
singular fibration could still exist, but it cannot arise as a global
effective torus action and must pass the map-level tests
\eqref{hrs8:eq:euler-fibres}, \eqref{hrs8:eq:van-kampen-fillings}, and
\eqref{hrs8:eq:mv-tests}.  The standard $T^4$ action gives an exact and
economical real $S^8$ multipole model, but its base is $B^4$, not $S^2$, and
its singular locus is a face stratification rather than three isolated
fibres.

The $3\times3$ $U(3)$ Gram construction gives an exact quotient boundary
$S^8$ and is the strongest precise realisation presently obtained for the
``one rung higher'' intuition.  Strict $SU(3)$ retains the determinant phase
and does not give the same quotient.  Keeping the original $T^4$ fibres and
placing $SU(3)$ on $X\setminus D$ is instead mathematically well typed: its
unframed bundle classes form $\mathbb Z^6$, every nonzero $c_2$ is an exact
nonextendable divisor defect, and global $SU(3)$ bundles on $S^6$ carry the
even-valued $c_3$ charge in \eqref{hrs8:eq:c3-clutching} while restricting
trivially to the regular locus.

None of these statements identifies a lepton, quark, or four-dimensional
Yang--Mills theory.

\paragraph{Sources.}
The almost-complex obstruction is cross-checked against A.~Borel and
J.-P.~Serre, \emph{Groupes de Lie et puissances r\'eduites de Steenrod},
Amer. J. Math. \textbf{75} (1953), 409--448, Proposition~15.1,
\texttt{doi:10.2307/2372495}, and P.~Konstantis and M.~Parton,
\emph{Almost complex structures on spheres}, Differential Geom. Appl.
\textbf{57} (2018), 10--22,
\texttt{doi:10.1016/j.difgeo.2017.10.010}.  The fixed-point identity used in
\eqref{hrs8:eq:torus-rank} is from S.~Kobayashi, \emph{Fixed points of
isometries}, Nagoya Math. J. \textbf{13} (1958), 63--68,
\texttt{doi:10.1017/S0027763000023497}.  For singular torus fibrations and
the coordinate sphere action, see respectively Y.~Matsumoto, \emph{Torus
fibrations over the 2-sphere with the simplest singular fibers}, J. Math.
Soc. Japan \textbf{37} (1985), 605--636,
\texttt{doi:10.2969/jmsj/03740605}, and H.~Ishida, Y.~Fukukawa, and M.~Masuda,
\emph{Topological toric manifolds}, Moscow Math. J. \textbf{13} (2013),
57--98, \texttt{doi:10.17323/1609-4514-2013-13-1-57-98}.  The theorem-level
calculations above are given in full and do not rely on these citations as
substitutes for omitted morphisms.
